\section{问题分析}
本题同时包含任务级离散决策、分钟级连续占用和小时级能源结算。若把任务直接压缩为小时平均负荷，会漏掉只在小时内短暂重叠却造成瞬时 GPU 峰值的任务；若把所有触及某小时的任务都按整小时计费，又会高估 AI 电量。因此本文先保留任务的真实起止区间，再用“是否正重叠”审计瞬时容量，用“实际重叠时长”核算 GPU-hour 和电量，最后把重建负荷传给能源模型。

\subsection{问题一：预测与末端调度分离}
预测对象是区域—任务类型—小时的到达 GPU 需求。训练、验证、测试必须按时间顺序切分，且多步预测时后续滞后项只能使用此前预测值，不能回填测试真值。日周期与周周期朴素模型提供可解释基线，Huber 回归处理尖峰需求并通过验证集选择超参数\cite{huber1964robust}。预测只用于精度评价；最后 24 小时调度仍读取第 2376 至 2399 小时实际到达的 538 个任务，从而避免“以预测量替代真实任务”的口径错误。

末端任务可能跨越第 2399 小时，因此调度时域延伸到第 2405 小时；第 2406 小时只用于能源和 SOC 结算。实时推理必须到达即开工，弹性任务可在到达与截止之间选择整数开工时刻。候选区域还需同时通过单向时延、GPU、IT 功率和设施功率下界筛选。

\subsection{问题二：五个方向的有限任务邻域}
全体 50000 个任务形成的区域—开工组合远超运行预算。本文不把“成本最小”“碳最小”等命名误写成数学最优，而是令五个名称分别表示搜索优先方向：成本、碳、新能源、服务和平衡。每个方向按任务类型、来源区域、到达时段、GPUh 分位和等待分位确定性分层轮转，每任务最多评价 8 个移动候选，总预算仍为 10000；五个方向采用不同目标函数与遍历偏移。增量数组只更新受影响的区域和小时，但每个方向结束后仍对全部任务执行一次独立全量审计。这样既扩大不同任务覆盖，又把搜索效率与最终可信性分开。

系统净成本定义为购电支出减售电收入。题面售电边界较大，因此负成本只表示这一核算边界内售电收入超过购电支出。为揭示结果来源，除总成本外还分别报告购电量、售电量、购电支出和售电收入，不以单一负数掩盖现金流结构。

\subsection{问题三：固定负荷下的离散 SOC 动态规划}
本问固定附件中的基准 AI IT 负荷，任务调度不再变化。每个区域设置 21 个等距 SOC 状态，并显式加入真实初始 SOC；逐小时只枚举满足容量、充放电功率、购售电互斥、能量守恒和终端 SOC 等式的转移。候选集合必须包含 no-action，另外按成本、碳、新能源、峰值、爬坡和平衡六种偏好形成策略。相同轨迹按哈希去重后先做候选集内 Pareto 支配筛选，再对所有不重复非支配方案统一进行 min--max 标准化，以成本、碳、新能源、峰值、爬坡权重 $(0.30,0.25,0.15,0.15,0.15)$ 计算到理想点的加权距离；不按“平衡”等名称优先\cite{miettinen1999nonlinear}。

高循环策略可能在未计退化成本时获得较好账面指标。故等效循环次数是解释变量而非新的硬约束：本文如实报告成本最小、平衡和新能源优先策略的循环水平，并用终端 SOC 放宽实验区分“终端库存约束影响”与“策略本身改善”。

\subsection{问题四：按情景动态重建邻域}
联合问题中，任务方案改变设施负荷，储能轨迹又改变任务移动的边际电价和边际碳信号。固定若干预制任务方案再逐情景挑选，会切断这种反馈。本文对每个情景从共同任务基线和真实初始 SOC 独立重启：先求当前负荷下的储能，再依据该储能形成的价格、碳、购电和弃电信号重新生成 600 个任务候选；接受改进后重建下一轮邻域。最多 8 轮，连续 2 轮无改善即停止。

17 个情景覆盖基准、4 个碳约束水平、3 个峰谷价差、3 个售电机制、3 个新能源水平，以及种子 2026、2027、2028 的新能源波动。所有情景保持任务集、时域、PUE、功率映射和评价公式一致。碳目标未达到只能表述为“声明启发式内未达到”，不能据此推断数学不可行。

\subsection{四问衔接与证据边界}
四问共享任务重叠、设施负荷、能源守恒和审计函数，但不直接拼接决策结果。问题二以问题一可行基线启动；问题三独立使用题定固定负荷；问题四再以问题二平衡任务方案为共同基线，并用与联合解相同的 21 状态网格和同一储能策略重算全部情景。输出可信性来自结构化结果、逐行 CSV、完整约束审计和两次隔离复算。由于不存在独立 Oracle，最高可信等级限定为 scenario-feasible。
